% --- I. Introduction ---
\section{Introduction}
An operating system (OS) has many jobs. It decides which program uses the processor, it shares the memory between programs, it sends requests to the disk, and it stops programs from using the same resource in a wrong way. In the labs of this course, we studied each of these jobs alone: a CPU scheduler in one week, a page-replacement algorithm in another week, and a disk scheduler later. In a real system, however, these mechanisms work at the same time and depend on each other~\cite{brief}. For example:
\begin{itemize}
    \item When a program needs a part of its memory that is not loaded (a \emph{page fault}), it must stop using the processor and wait for the disk.
    \item When the disk finishes, the program becomes ready to run again, and the CPU scheduler must decide when it runs.
    \item When a program waits for a lock that another program holds, the scheduler cannot choose it until the lock is free.
\end{itemize}
So a slow disk can make the CPU idle, and a small memory can make the whole system slow. These effects are hard to see when each mechanism is studied alone.

\subsection{Goal of the Project}
The goal of this project is to build one simulator, called \ossim, in which all of these mechanisms run together, and then to use it as a tool for experiments. A \emph{simulator} is a program that imitates the behavior of another system. \ossim{} does not run real programs. Instead, it reads a text file (a \emph{workload}) that describes a set of imaginary processes: when they arrive, how long they use the CPU, when they do input/output (I/O), which pages of memory they use, and which locks they take. Then it simulates what an OS would do with them and measures the result.

The command line of \ossim{} selects the policies to use: the CPU scheduler, the page-replacement policy, the disk scheduler, the number of memory frames and the disk parameters. In this way, the same workload can be run with different policies and the results can be compared fairly.

\subsection{What the Team Built}
The course gave a starting version of \ossim{} (the \emph{scaffold}). It already had the main loop of the simulator, the simplest policies (FCFS and Round Robin for the CPU, FIFO for memory and FCFS for the disk) and the code for semaphores and resources. Our team completed the rest in two checkpoints:
\begin{itemize}
    \item \textbf{Checkpoint 1} (processes, scheduling and synchronization): four more CPU schedulers (SJF, SRTF, Priority and MLFQ) and deadlock detection.
    \item \textbf{Checkpoint 2} (memory and storage): two more page-replacement policies (LRU and Clock), three more disk schedulers (SSTF, SCAN and C-LOOK) and a workload generator.
\end{itemize}
For every policy, we first solved an example by hand and then checked that the simulator gives exactly the same answer. After that, we used the simulator for experiments, and we extended it with one extra mechanism (the \emph{theme}).

\subsection{Structure of this Report}
Section~\ref{sec:background} explains the OS concepts that the rest of the report uses. Section~\ref{sec:design} describes how the simulator is built. Section~\ref{sec:validation} shows how we checked it against hand-computed examples. Section~\ref{sec:experiments} presents two experiments, and Section~\ref{sec:theme} presents the theme and its experiment. Section~\ref{sec:conclusion} summarizes the results. All the code is available in the team repository~\cite{repo}.
